KernelBench v3 results reinforce two observations. First, baseline-beating kernels are rare even for top models. This aligns with METR's finding that kernel benchmarks can overstate capability when task design is weak~\cite{metr2025}. Second, performance collapses on Level 4 tasks, which emphasize modern attention and quantization patterns~\cite{fp8-formats,int4-quant}. These tasks are not just harder because of novelty but also because some baselines are already optimized Triton implementations~\cite{triton}, shrinking the remaining optimization headroom.

We also observe that correctness is not the limiting factor for all models. Several models achieve respectable correctness rates but lower baseline-beating rates, suggesting that the remaining gap is performance engineering rather than syntax or functional accuracy. For practical kernel engineering, models must reason about memory layout, kernel fusion, and hardware-specific bottlenecks.

Finally, cost remains a key practical constraint. The benchmark logs per-run token counts and estimated costs, enabling analysis of accuracy and speedup as a function of spend. This makes KernelBench v3 useful for evaluating cost-performance tradeoffs, not just raw capability.
